\chapter{Conclusion}\label{ch:conclusion}

The ODMI runs every year over 36 countries and 5,148 question-country pairs, each answer checked by hand \parencite{edp2025}, and models now are able to match humans at far lower cost on comparable text \parencite{heseltine2024,tornberg2025}. \textcite{davis2012} argue that indicators like the ODMI shape how states can be ranked and how resources will be allocated, and that the authority for this rests on contestability, where a country is able to see and dispute the evidence held against its score. Automation threatens contestability directly, because a generative assessor can produce a fluent-sounding answer without an adequate basis and the `black-box' opacity of these LLMs makes results difficult to contest. \textcite{fumega2026} ask whether AI can take over indicator evaluation without compromising contextual nuance or democratic legitimacy. The basis of evidence for responses is more important for contestability than the answers themselves, and automation needs to accommodate this.

Section~\ref{sec:criteria} assembled a novel framework for using generative models in evaluating assessments of this kind, taking inspiration from similar works for deterministic and predictive models \parencite{yung2022,williamson2012,schenk2024}. These provide a measurable way to ensure contestability: Convergent Validity, Generalisability, Subgroup Equity, Attributability, Selectivity and Reproducibility. These set out what an automated assessor has to show for its results to stay contestable. The swarm was built against them, with an adversarial Verifier to challenge each answer and an abstention channel to return nothing when the evidence is insufficient.

The swarm ran over eight countries, a total of 1,144 question-country pairs, committing an answer on 636 and declining 508, and matching the published 2025 key on 70.1\% of the committed answers that could be scored. The open-world assumption of Section~\ref{sec:open-world} governs that result. A country reporting on itself works from a complete record where silence means absence \parencite{reiter1978}, and an agent searching the open web cannot be certain. Every search returned material and no abstention came from an empty result set, so the system never fabricated evidence. What it could not do was find evidence that a practice is absent. It matched 87.0\% of the yes golds it committed on and 48.9\% of the noes, and scored 17 points lower than if it arbitrarily answered ``yes'' to each question.

Swapping the Verifier from adversarial to corroborative over the same 1,144 pairs had no significant effect on accuracy or coverage. The Researcher and the Verifier are the same Claude model separated by a prompt, and changing the prompt made no difference. What did move the result was adding calls, with the Verifier lifting coverage from 26.6\% to 46.0\% and the Adjudicator to 55.6\%, each buying answers and not precision. Current research is locating the benefit elsewhere. \textcite{chowdhury2026} assign a different model to each of their judges and put 7.5 of their 10.0-point gain on progressive evidence retrieval, and \textcite{zhu2026} attribute the benefit of debate to agent diversity. Both locate the gain somewhere other than the Verifier's stance, in fresh retrieval and in model diversity. This design retrieved fresh evidence on every retry but ran one model throughout.

Attributability, Reproducibility and Convergent Validity are mostly met. Every committed answer carried a passage found in the text the Researcher read, so the fabrication issue raised in Attributability's first condition is resolved. The Verifier is the reasonable reader the second condition asks for and halted 208 of the Researcher's claims as off-target, though an external model or human reviewer would have been a stronger test to catch the over-reads that were let through. Reproducibility holds for the most part. Receipt-taking at every step allows for easy replication of any run at any stage. The stability condition is a mixed picture. There is an insignificant 1.4 points difference in commit-accuracy over three runs, and a high 92\% agreement rate on labels it committed to, however answers near the confidence threshold were not consistent, leading to a 30\% variation on what questions were committed. For Convergent Validity, there is a 70.1\% agreement rate between the swarm's answers and the human responses. The independent recompute of deterministic questions diverged from the answer key 50\% of the time, though only four of the eight countries were computable.

Selectivity, Subgroup Equity and Generalisability all fail for the same reason. The confidence that governs the decision to answer measures how well a claim is evidenced. This works well for the positive class of answers but leaves the negative class consistently below the threshold required to commit. Selectivity fails on this directly, since its second condition asks the stated confidence be well-calibrated, and on the negative class, calibration error runs at 0.299 against 0.089 on the positive class. False positives are a consequence of this misalignment. Where little has been published about a question, a negative answer falls below the threshold and the system declines. However, where a large quantity has been published around the question, which while relevant does not directly answer the question, the material sitting near the question lifts it over the threshold and the system commits an incorrect ``yes''. The swarm declines 61.5\% of Bosnia and Herzegovina's pairs and has a false-positive rate of 12.8\% on its negative golds. For Finland, it declines 25.9\% of the questions with a false-positive rate of 48.3\%. This is Subgroup Equity failing. Furthermore, the Quality dimension has a commitment rate of 42.7\%, at an 18.8\% false-positive rate against Policy's 74.2\% and 50.0\%, which is Generalisability failing.

\textcite{jacobs2021} identify that in assessments like the ODMI, the construct of the quality of openness can diverge from its operationalisation, the assessment. In this case, a team that says it monitors its search terms will satisfy the question without actually having produced meaningful and verifiable evidence. Given this, the questionnaire in its current state could be seen to be flawed. Of the 143 ODMI questions about 20\% ask a team about its own internal practice where the supporting evidence can be seen as unverifiable. Moreover, 15.1\% of the 5,148 answers were returned with no evidence at all, and the human verifiers at Capgemini alter 10\% of answers submitted by the national teams.

There are considerable opportunities for the ODMI itself to be reformulated to address the weaknesses outlined here, so that every answer becomes publicly evidenceable. This would open far more of it to automation, as an automated assessor reading the public record doesn't rely on the self-reporting the ODMI currently depends on. This would also make the ODMI a better test of openness, because it would be searching for evidence of each country's practice rather than how they self-report, which narrows the gap between the construct of openness and the operationalisation of the assessment. The European Commission treats open data as infrastructure for wider digital policy and for artificial intelligence in particular \parencite{edp2025}. Therefore, improving the ODMI as a test of openness would enable countries to be assessed more accurately for the honesty of their practices, and any shortcomings would be reflected in their results.
